% Single slide: Reinforcement Learning + PPO (for deck insertion).
% Compile: pdflatex E90_slide_ReinforcementLearning_and_PPO.tex
% Or copy the frame into a larger Beamer presentation.
\documentclass[aspectratio=169]{beamer}
\usetheme{default}
\usecolortheme{default}
\setbeamertemplate{navigation symbols}{}
\setbeamertemplate{footline}[frame number]

\title[E90]{E90 Mid-Semester Check-in}
\author{Howard \& Matthew}
\date{March 2026}

\begin{document}

\begin{frame}{Reinforcement Learning \& PPO}
  \footnotesize
  \textbf{Reinforcement learning}
  \begin{itemize}\setlength{\itemsep}{0.15em}
    \item \textbf{Agent} vs.\ \textbf{environment} (sim robot + terrain)
    \item Loop: \textbf{obs} $\rightarrow$ \textbf{action} $\rightarrow$ \textbf{reward}
    \item \textbf{Policy}: obs $\rightarrow$ action
    \item Maximize \textbf{long-term return}; learn by \textbf{trial and error}
    \item No fixed joint script---reward tells us what ``good'' means
  \end{itemize}
  \vspace{0.35em}
  \textbf{PPO (Proximal Policy Optimization)}
  \begin{itemize}\setlength{\itemsep}{0.15em}
    \item \textbf{Policy-gradient} algorithm (updates from rollout data)
    \item \textbf{Actor} = policy; \textbf{critic} = value network
    \item \textbf{Clipped} updates $\Rightarrow$ small, stable policy changes
    \item Strong fit for \textbf{continuous} actions (e.g.\ 12 joints)
  \end{itemize}
\end{frame}

\end{document}
